\clearpage
\fancyhead[L]{\small PIAZZI. LEZIONI DI ASTRONOMIA.}
\hypertarget{native-622-lower}{}
\begin{center}\small Göttingische gelehrte Anzeigen. No. 4. Pages 33..37. 6 January 1820.\\[3pt]\rule{.72\linewidth}{.3pt}\end{center}

\noindent\emph{Lezioni elementari di astronomia ad uso del real osservatorio di Palermo. Palermo {\upshape 1817}. Dalla stamperia reale.} Volume I, 240 pages. Volume II, 446 pages. Small quarto.

\begingroup\brokenpenalty=10000\relax
In the preface, the author of this textbook, the renowned and distinguished \textsc{Piazzi}, explains that the book grew out of the lectures he had given at the Palermo Observatory and had not originally been intended for the public. He adds that the authorities responsible for public instruction ordered its publication at the students’ urgent request. One expects no substantive advance in the science itself from writings of this kind. The main thing is a clear, thorough organization of the exposition in keeping with its purpose. What the exposition requires in turn varies greatly with the students’ prior knowledge and aims. The requirements for mere amateurs differ from those for students trained in mathematics. They differ again for someone with the sensibility to revere science as one of the most perfect and noblest products of the human \emph{mind}, who takes nothing on trust but demands strict logical coherence throughout, and who prefers the finest exposition even if it requires more than the usual effort and ancillary knowledge. Finally, they differ again for someone who pursues only the lower aim of becoming proficient in observing and calculating celestial phenomena, to the extent that more limited mathematical training permits. The author does not explicitly say which class of students he actually intends the book for, but he does seem to have had chiefly the last class in mind. From an astronomer like \textsc{Piazzi}, who has rendered such great service to several branches of astronomy, one could moreover expect that even an elementary work of this kind was bound to contain many distinctive views, noteworthy remarks, and judgments. In this respect it is not without interest even for astronomers.

The author has divided his work into seven books. The first gives a general survey of celestial phenomena as far as rough observation alone can take us, without instruments or special mathematical knowledge. It therefore deals with the general diurnal motion, the points and circles adopted on the celestial sphere, the shape of the Earth, the distinction between fixed stars and planets, the Sun, the Moon, the planets and comets, the ecliptic and the other circles of the celestial sphere determined by it, the zodiac and constellations, the precession of the equinoxes, eclipses, and the first inequalities of the planets and their explanation by the Copernican system. The second book deals with the knowledge needed for the new astronomy. Here we find some material on solving spherical triangles and an overview of the principal astronomical instruments now in common use, such as telescopes, a plumb line, a spirit level, verniers, micrometers, a transit telescope, a clock, a full circle, a quadrant, a zenith sector, an equatorial, and repeating circles (\textsc{Piazzi} judges the last of these unfavorably, though he does seem to have somewhat exaggerated what can be said \SourceLeaf{623}against them and has indeed been said by others. Had he himself been required to observe with an instrument of this kind, says \textsc{Piazzi}, he would have accomplished \emph{not even a hundredth} of what he did). The book then deals with refraction, parallax, aberration and nutation insofar as the observations need to be corrected for them. It goes on to discuss time and the methods of determining it, and finally the observations used to determine the altitude of the pole, the obliquity of the ecliptic and the positions of celestial bodies relative to the equator and the ecliptic. The third book then deals exclusively with the fixed stars. It treats the calculation of precession, nutation and aberration in great detail, then considers the annual parallax of the fixed stars, their proper motion, the star catalogues and the peculiar features of the fixed stars. We cannot pass over a peculiar blunder in the determination of the magnitude of the annual precession. \textsc{Piazzi} subtracts from the observed annual change in the longitude of 15 fixed stars, which follows from comparing \textsc{Bradley}'s determination for 1755 with \textsc{Piazzi}'s own for 1805, both the effect of the displacement of the ecliptic and that of \emph{proper motion}. He then takes the mean of all the results (which, incidentally, vary greatly). This is plainly circular reasoning, since the proper motion of each star can be ascertained only by comparing the observed motion with the precession already assumed to be known. In a consistent calculation all the stars should have yielded one and the same result, namely the very result taken as the starting point for finding the proper motions.

The subject of the fourth book is the elliptical motion of the planets, drawn for the most part from the Theoria motus corporum coelestium, with all the proofs omitted, however.

The fifth book considers the celestial bodies of our solar system individually. It is remarkable how \textsc{Piazzi} discusses, on p. 182, the late \textsc{Schröter}'s determination of the rotation period of Venus. `Niente di meno non manca, chi non poco dubiti di tali scoperte, e le consideri come non molto dissimili dall' altra del satellite, che \textsc{Short} e qualche altro credeva di aver veduto intorno a quaesta pianeta. Una illusione ottica, un fenomeno non bene osservato, accompagnati da un poco di immaginazione, hanno talora fatto avanzare delle congetture, e sognare delle esseri, che ben presto son poi caduti in piena dimenticanza (Yet \textsc{Piazzi} seems to have forgotten here that on p. 46 of the first volume he himself had declared the existence of a satellite of Venus highly probable, a satellite in which hardly any other astronomer still believes). Likewise, on p. 230 he takes issue with the burning lunar volcanoes that some astronomers claimed to have seen.

The sixth book deals with eclipses, occultations of stars and transits of the inferior planets across the solar disc, and the seventh finally deals with comets. For determining their orbits, only the indirect method is given here, in the same extremely imperfect and, one might say, haphazard form that astronomers found to their astonishment in a well-known astronomy textbook published a few years ago.
\par\clearpage\endgroup
